\chapter{结论与展望}
\label{sec:conclusion}

\section{研究结论}

本文围绕考虑顾客选择的服务网络设计问题，提出了行程聚合与动态细化框架（\algoNotation）。主要工作和结论概括如下。

在建模方面，本文建立了考虑顾客选择的服务网络设计混合整数规划模型，在时空网络上同时处理服务调度、资源分配和顾客需求分配决策，并通过基于GAM/SBLP的线性选择表示将顾客选择行为纳入约束体系。

在行程聚合方面，本文利用当前离散化诱导的路径签名将原始行程划分为等价类，每个等价类由一条超级行程代表。超级行程的吸引力被构造为组内可兼容原始行程组合吸引力总和的上界，从而在压缩行程变量规模的同时避免了时间聚合带来的虚假收入。

在求解框架方面，本文基于超级行程的构造，得到了聚合下界模型、完整网络上界修复模型以及松弛解可实现性检验机制，形成完整的双侧界求解框架。下界模型在聚合网络和超级行程空间上求解，给出原问题最优值的有效下界；上界修复模型利用频率固定和超级行程收入约束在完整网络上生成可行解；可实现性检验通过识别不可能出发、间隔出发违规和虚假收入三类非法情形，指导时间点的定向添加。在不受最大迭代次数截断时，算法在有限步内收敛到全局最优解。

在实验验证方面，本文在真实航空排班实例上测试了框架的计算效果。五个实例的可行行程数从约48万到142万不等，\algoNotation 均在5小时内达到$<1\%$的最优性间隙。作为对比，直接求解完整模型在中小规模实例上仍保留约$7\%$--$14\%$的间隙，在百万级行程变量实例上则未能给出可用结果。

\section{研究不足与展望}

本文的工作仍有若干局限，以下几个方向值得进一步探索。

实验验证集中在航空排班场景。尽管建模框架和求解方法本身不依赖于特定交通方式，但尚未在铁路运行图编制、城际客运班次规划或城市公交网络设计等场景中进行系统测试。后续研究可以开展更加广泛的跨领域验证，以评估方法在不同服务网络设计问题中的适用性和性能表现。

算法效率方面仍有改进余地。例如，可利用相邻迭代间的解信息设计热启动策略，或结合机器学习等数据驱动方法预测有潜力的行程进行智能聚合和划分，以进一步提升求解效率。此外，算法的并行化实现也是一个值得探索的方向，利用并行技术同时求解多种行程聚合方案，可能获得更强的鲁棒性和更好的性能表现。

顾客选择行为采用GAM/SBLP框架建模，尚未考虑更复杂的选择结构（如嵌套Logit、混合Logit）或更复杂的需求动态（如时变需求、顾客忠诚度）。未来研究可探索更丰富的选择模型，以更准确地捕捉顾客行为特征。不过，这也将带来更大的建模和求解挑战，混合整数规划的线性表示可能不再适用，需要开发新的建模和求解方法，如基于非线性规划或启发式算法的框架。这些方法需要在保持求解效率的同时，能够处理更复杂的选择结构和需求动态，是很有挑战但也很有价值的研究方向。
